\section{Method}
\label{sec:method}

We analyze the local copy of GitSkills with Python and Polars DataFrames.
Each notebook sets one shared random seed when it loads the tables, and every sampled draw we report uses a fixed seed, so a rerun returns the same rows.

\subsection{Study population}
\label{sec:population}

To keep the population small and well defined, a skill stays in the study only when every rule below holds, whatever its body length.

\begin{itemize}
  \item The text is counted once, so copies of the same skill are left out (\textit{dedup\_primary}).
  \item The header metadata at the top of the file is valid (\textit{frontmatter\_valid}).
  \item The file is named \textit{SKILL.md} or \textit{skill.md}, and its body is plain text rather than JSON or CSS.
  \item The file sits in the usual skill folder under \textit{.claude/skills} or in another skills directory (\textit{location\_class}).
\end{itemize}

The data checks in Section~\ref{sec:results} used the full tables, so these rules were not applied to them.

\subsection{Taught-tool extraction}
\label{sec:extraction}

Extraction reads only the skill body, never the companion text.
We extract names with GLiNER~\cite{zaratiana2023glinergeneralistmodelnamed}, which scores each hit, and a cutoff of~$0.8$ on that score decides which hits we keep.
The cutoff and a target accuracy of~$80\%$ are both arbitrary choices for this paper, and the manual validation in Section~\ref{sec:validation} checks the cutoff against that target.
In the label-only run in Section~\ref{sec:results}, labels alone accepted false positives, and Section~\ref{sec:relation} proposes improvements to this step.

\subsection{Manual validation}
\label{sec:validation}

Every step that relies on a model will go through a human quality gate.
Those steps are taught-tool extraction, the choice of an owning repository when a tool is not already linked, and license gaps when a model helps fill them.

We will label the taught tools in a seeded random sample of the population, sized before anyone sees the model output, and report precision and recall against those labels.
Precision is the share of extracted tools that are truly taught tools, and recall is the share of truly taught tools the model finds.
Two authors will label a subset independently, and we will report their agreement.
A stage that falls short will run again with new labels or a new cutoff.
Once it meets the target accuracy, we will fill the remaining gaps manually and merge the automatic and manual results, dropping entities that do not exist, before the next stage begins.
Repository resolution and model-assisted license steps will use the same gate on their own samples, as Section~\ref{sec:future} describes.

Extracting every tool manually, or checking every model output, would take about the same effort, and either path would limit how many skills, tools, and repositories we can cover.
We will therefore accept a bounded residual error under the sample gate so the study can scale, in the same way a survey estimates from a sample instead of visiting every case.
In extraction, the gate is meant to catch a model that marks the hosting folder, the publishing repository, an attribution link, the skill filename, or GitHub inside an address as a taught tool.
